\section{Ketertelusuran Rancangan Penelitian}

Setelah variabel, rancangan algoritma, konstruksi dataset, desain eksperimen, metrik, dan teknik analisis ditetapkan, Tabel \ref{tab:pemetaan-kebutuhan} merangkum ketelusuran dari akar masalah, kebutuhan artefak, keputusan rancangan, hingga bukti evaluasi. 
Matriks ini memastikan setiap komponen CoDA-EDA memiliki dasar permasalahan yang jelas dan setiap klaim penelitian memiliki mekanisme pengujian yang sesuai.

{\small % <--- PERINTAH PENGECILAN FONT DITAMBAHKAN DI SINI
\begin{longtable}{@{} >{\centering\arraybackslash}p{1cm} >{\raggedright\arraybackslash}p{3cm} >{\raggedright\arraybackslash}p{3cm} >{\raggedright\arraybackslash}p{3.0cm} >{\raggedright\arraybackslash}p{3cm} @{}}
\caption{Desain artefak dan pemetaan kebutuhan}
\label{tab:pemetaan-kebutuhan} \\
\toprule
ID & Masalah / akar penyebab & Kebutuhan artefak & Keputusan desain & Bukti evaluasi \\
\midrule
\endfirsthead
\caption{Desain artefak dan pemetaan kebutuhan (lanjutan)} \\
ID & Masalah / akar penyebab & Kebutuhan artefak & Keputusan desain & Bukti evaluasi \\
\midrule
\endhead
\bottomrule
\endfoot
M1 & Populasi eksplisit membebani memori & Memori tidak bergantung pada $NP$ & PV tanpa populasi eksplisit & Kurva memori vs $NP$; analisis \& regresi vs $D$ (Gambar IV.x) \\
M2 & Model univariat mengabaikan dependensi & Menangkap dependensi dengan biaya terbatas & Pohon Chow-Liu orde satu ($O(D)$ via Prim) & Perbandingan terhadap baseline compact univariat (compact-GA/PSO); kualitas per dimensi/korelasi (H1) \\
M3 & Pembaruan struktur berulang mahal & Struktur diperbarui hanya saat diperlukan & Pemisahan pembaruan parameter--struktur + deteksi drift (JSD, $\tau$) & Ablasi mekanisme pembaruan (warm vs always); trigger rate; overhead deteksi (H3) \\
M4 & Kebijakan berubah setelah drift & Re-optimasi tanpa mining penuh & Warm-start dari model sebelumnya & Ablasi warm-start (warm vs cold); waktu menuju target kualitas + uji non-inferioritas (H4) \\
M5 & Benchmark berpotensi tidak adil & Representasi, fitness, anggaran, operator setara & Kontrol eksperimen bersama (kesetaraan implementasi) & Tabel konfigurasi, sanity check, kontrol kesetaraan \\
\end{longtable}
} % <--- PENUTUP GRUP

Dengan demikian, alur argumentasi dari Bab I hingga Bab IV dapat ditelusuri secara konsisten. 
Setiap masalah pada konteks \textit{edge-IoT} dihubungkan dengan mekanisme rancangan yang menjawabnya serta eksperimen yang digunakan untuk mengevaluasinya.